\usepackage[
	backend=biber,
	style=ieee
]{biblatex}

\addbibresource{references.bib}

\usepackage[
	backend=biber,
	style=ieee
]{biblatex}

\paragraph{
	BLDCs have largely replaced DC motors for robot joint actuation \cite{limb_exoskeleton} due to __ and  __. Modern robot joint predominantly assume BLDC-driven joints. BLDC motors differ from brushed DC motors in that they are driver by three phase currents that energize their stator windings. By controlling the magnitude and timing of these currents in the correct sequence, the controller creates a rotatiang magnetic field that causes the permanent-magnet rotor to rotate. Providing the correct voltage at the three phases in the right sequence is the job of the Motor Controller.

}

\paragraph{
	A member of the family of sinusoidal motor controllers is the sensored Field Oriented Controller (FOC) \cite{Wishart1990FOC}, which works by transforming th phase currents from the Iu, Iv, Iw frame to a 2D voltage vector, the alpha beta frame, oriented wrt to the stationary stator. It represents the strenght of magnetic field at the three stator windings as a vector aligned with the direction of the magnetic field of the stator windings. It then transforms this AlphaBeta ($I_{\alpha, \beta}$) frame into the DirectQuadrature frame, which is vector that orients the magnetic fields wrt the rotating rotor. The D represents the current energizing the magnetic field that causes the rotor to point towards the incumbent stator winding, therefore, its wasted current, as we would like forward rotation. The Q represents the current energizing the stator windings that will cause forward rotation. To calculate I_DQ, we need the current orientation of the rotor. A rotary encoder/sensor is used to get this orientation (hence "sensord" FOC).

	After transforming the phase currents u,v,w to the DirectQuadrature frame, we can now solve for the commanded torque by implementing two current PI control loops.

	one regulates d-axis to I^*_{d} (0) and other rgulates q-axis to , and one regulating q-axis term to desired torque $I^*_q$, producing V_d and V_q respectively, which are voltage commands generated by current controllers to regulate d-axis and q-axis currents to their desired values, $I*_d and I*_q$. V_d and V_q are transformed to the stationary, stotor oriented voltage vector V_ab using inverse park.

	A 3-phase inverter can produce any of the size non-zero voltages:
	000,001,010,011,100,101,110,111, of which 001-101 form a hexagon. 100 for example means phase u is +, phases v and w are -.

	and then using SVPWM converted to 3 inverter duty cycles, d_u, d_v, d_w, sent to the inverter, which converts the duty cycle to voltages V_u, V_v, V_w.

	The benefits of this control architecture is manifold:
	Smoother torque ripple, near constant torque, important for robots that require precise and repeatable motion. They produce desired torque with minimal wasted currnt. It provides precise current control, and since torque is proportional to currnet, enables responsibe and accurate torque control. Excellent low speed control, remains smooth even near 0 speed, better than trapezoidal which produces discrete torque steps (cogging and vibration).
}

\paragraph{
	another technique is feed forward compensation \cite{Heydari2010FOCSVPWM} . a basica currnt PI controller measures current error $e = I_{ref} - I_{measured}$ and regulates that, overcoming motor resistance, inductance and back-EMF. Feed forward accounts for _ in its model and automatically compensates them. Voltage equations are $V_d &= R_s I_d + L_d \frac{dI_d}{dt} - \omega_e L_q I_q$, $V_q &= R_s I_q + L_q \frac{dI_q}{dt} + \omega_e L_d I_d + \omega_e \lambda_m$. This makes contorller faster as it doesn't fight motor dynamics. e
}

\paragraph {
	An issue with Sinusoidal PWM is that it doesnt maximise DC bus voltage utilization. SVWPM is a modulation technique that solves that by extending the linear modulation range by ~15 percent compared to sinusoidal PWM.
	There are different variants: continous SVPWM, discontinous PWM (DPWM0, DWPM1, DPWMmin DPWMmax), overmodulation SVPWM, bus-clampd SVPWM.
}


\paragraph{
	Benefits of Odrive
	Cascaded control loops: a fast current/torque loop inside a slower velocity loop, inside a slower position loop. Torque mode bypasses the outer loops, while velocity and position modes reuse the same inner current controller. ODrive Documentation

	FOC current control: phase currents are transformed into \(I_d\) and \(I_q\), regulated using PI controllers, then converted back into inverter duty cycles.

	PWM-synchronized current sampling: ADC measurements should occur at carefully selected points in the PWM cycle, rather than at arbitrary times, to reduce switching noise and ensure valid current readings.

	Hardware PWM generation: STM32 timers generate complementary three-phase PWM signals, typically with center-aligned PWM and inserted dead time. The CPU updates duty-cycle registers rather than manually toggling pins.

	Direct peripheral triggering: timers trigger ADC conversions, and ADC completion triggers the control interrupt. This minimizes timing jitter and unnecessary software scheduling.

	Sensored rotor feedback: ODrive supports incremental encoders, Hall sensors, SPI absolute encoders, RS485 encoders and other feedback devices, depending on the model. ODrive Documentation

	Separate commutation and load encoders: one encoder can measure the motor rotor for FOC, while another measures the output after a gearbox. This is useful for compensating gearbox backlash and compliance. ODrive Documentation

	Encoder interpolation and estimation: raw encoder readings are filtered or estimated to obtain smoother electrical angle and velocity signals between encoder edges.

	Automatic encoder-offset calibration: the firmware determines the relationship between encoder angle, motor electrical angle, phase order and direction. ODrive Documentation

	Automatic motor-parameter identification: ODrive measures phase resistance by applying controlled DC current and estimates inductance from current ripple under an alternating voltage. ODrive Documentation

	Parameter-based current-loop tuning: measured motor resistance and inductance can be used to calculate appropriate current-controller gains rather than relying entirely on manual tuning.

	Torque, velocity and position command modes: the same actuator can accept direct torque commands, velocity commands or position commands. ODrive Documentation

	Feedforward commands: position control can include velocity and torque feedforward, while velocity control can include torque feedforward. This improves tracking because feedback does not have to generate the entire command from error alone. ODrive Documentation

	Trajectory generation: position setpoints can pass through trapezoidal trajectory generation, velocity ramps or input filtering instead of being applied as abrupt steps. ODrive Documentation

	Anticogging compensation: the controller can learn a position-dependent cogging-torque map and add compensating torque during operation. ODrive Documentation

	Current and torque limits: software limits prevent commands from requesting excessive phase current, while harder limits can immediately disarm the inverter.

	Velocity and position limits: outer control loops constrain mechanical motion separately from electrical current limits.

	Bus-voltage monitoring: the controller monitors the DC-link voltage for undervoltage and overvoltage conditions.

	Regenerative-energy handling: braking energy can be returned to the supply or dissipated through a controlled brake resistor or regeneration clamp. ODrive Documentation

	Thermal monitoring and derating: MOSFET and motor thermistors can reduce allowable current before temperatures become unsafe, instead of only shutting down after overheating. ODrive Documentation

	Fault-latched state machine: calibration, idle, closed-loop control and fault states are explicit. Dangerous faults disarm the PWM outputs rather than relying on normal control logic.

	Watchdog protection: the actuator can automatically stop if commands or heartbeat messages from the host disappear.

	Hardware gate-driver protection: overcurrent, shoot-through and gate-driver faults should be handled by hardware where possible, because firmware may react too slowly.

	Dead-time and complementary-output management: high-side and low-side MOSFETs are prevented from conducting simultaneously.

	Deadline monitoring: the firmware detects when the real-time control computation does not finish before the next PWM cycle.

	Separation of fast and slow tasks: current control and encoder sampling run in the hard real-time path; communication, configuration, logging and calibration management run at lower priority.

	Non-blocking communication: USB, CAN and UART processing should not delay the motor-control interrupt.

	Multiple command interfaces: ODrive exposes CAN, USB, UART, PWM, analog input and step/direction interfaces. ODrive Documentation

	Structured object-based configuration: motor, encoder, controller and communication settings are organized into independently configurable subsystems and stored persistently.

	Runtime telemetry: currents, position, velocity, temperature, bus voltage, errors and controller states can be inspected for tuning and fault diagnosis.

	Explicit unit conventions: commands use physical quantities such as amperes, torque, turns, turns per second and volts rather than raw PWM counts.

	Modular hardware abstraction: control algorithms are separated from encoder drivers, communication drivers and board-specific peripherals, making the architecture easier to extend.

	Field weakening (FW). Allows operation beyond the motor's base speed by commanding a negative \(I_d\). ODrive historically did not implement this because it prioritized other features. ODrive Community

	Maximum Torque Per Ampere (MTPA). Instead of always using \(I_d=0\), compute the optimal \((I_d,I_q)\) pair that produces the required torque with minimum current. This improves efficiency, especially for interior PMSMs. arXiv

	Model-based feedforward. Add back-EMF and cross-coupling compensation to reduce the burden on the PI controllers. ODrive already has some support, but more complete implementations exist in the literature. arXiv

	Disturbance Observer (DOB). Estimate unknown load torque and compensate for it. This greatly improves disturbance rejection and is widely used in precision robotics. arXiv

	Online parameter identification. Continuously estimate resistance and inductance as the motor heats up instead of assuming constant values.

	Adaptive PI gains. Schedule current-controller gains based on operating point or identified motor parameters.

	Model Predictive Current Control (MPCC/MPC). Instead of PI + SVPWM, predict the effect of each switching state one step ahead and choose the optimal one. It gives excellent dynamic response but requires more computation. DOI

	Dead-time compensation. Correct inverter nonlinearities caused by MOSFET dead time, improving low-speed current accuracy.

	High-frequency injection. Enables sensorless operation at zero speed by exploiting rotor saliency.

	Current FOC + SVPWM (baseline)
	Feedforward decoupling
	Dead-time compensation
	Online motor parameter estimation
	Field weakening
	MTPA
	Disturbance observer
	Model predictive current control (if enough MCU performance is available)
}
